\documentclass[10pt,a4paper]{amsart}
\usepackage[margin=19mm]{geometry}
\usepackage{amsmath,amssymb,mathtools}
\usepackage[scr=rsfs,cal=euler]{mathalfa}
\usepackage{xfrac}
\usepackage[unicode,hidelinks,bookmarks=false]{hyperref}
\newcommand{\ie}{i.e.\ }
\newcommand{\cf}{cf.\ }
\newcommand{\resp}{respectively\ }
\newcommand{\Subsection}{\subsection}
\providecommand{\nobreakdash}{\nobreak\hbox{-}}
\setlength{\emergencystretch}{2em}
\multlinegap0pt
\usepackage{bm}
\usepackage{bbm}
\usepackage{dsfont}
\usepackage{xspace}
\usepackage{cancel}
\usepackage{mathtools}
\usepackage{amsthm}
\makeatletter
\@ifundefined{theorem}{\newtheorem{theorem}{Theorem}}{}
\makeatother
\title[Parallel Machine Scheduling with a Singler Server and Loadin]{Parallel Machine Scheduling with a Singler Server and Loading-Unloading Operations}
\author{Keramat Hasani, Frank Werner}
\date{Source: arXiv:2608.12087v1 (CC BY 4.0)}
\begin{document}
\maketitle

\noindent\emph{Source and licence.} Parallel Machine Scheduling with a Singler Server and Loading-Unloading Operations. Version arXiv:2608.12087v1 (CC BY 4.0), distributed under the Creative Commons Attribution 4.0 International licence, as recorded in the arXiv licence metadata for this exact version and shown on the abstract page https://arxiv.org/abs/2608.12087v1. Full rights notice: licenses/arxiv-2608.12087-CC-BY-4.0.txt.\par\smallskip

\noindent\textit{Editorial scope.} Counted unit: the Theorem whose source label is \texttt{thm:variable-m-strongly-np-complete} in the article (source0.tex lines 321--328, source label \texttt{thm:variable-m-strongly-np-complete}), delivered together with the article's own complete original proof of it (source0.tex lines 330--506, one \texttt{proof} environment of 177 lines). No mathematics is written, repaired or re-derived: statement and proof are the article's own text line for line. Required context retained uncounted: none. Every definition, display and label of those lines is retained unchanged. Omitted from this package and not redistributed: 1--320, 329--329, 507--1818. Nothing inside a retained excerpt is omitted or altered: every retained body is delivered exactly as the article's lines are written, so no omission receipt of any kind accompanies this package. Citations: the retained excerpts carry no citation command at all, so no bibliography entry of the article is transcribed and no reference is redistributed. Reference adaptation: none was needed and none was made; every reference inside the delivered unit resolves inside it, so no unresolved reference or citation is produced. Printed slips kept literal: no printed word, symbol, hypothesis or number of the counted statement or of its proof was altered. Layout and class: the article's own class is \texttt{article} with packages including geometry, amsmath, amssymb, amsthm, booktabs, graphicx, hyperref, tikz, float; the wrapper is the lane's pinned \texttt{amsart} editorial wrapper at 10pt on A4, which restates the theorem-like environments the retained text needs and adds the article's own macro definitions that the retained text needs (none), with extra wrapper packages (bm, bbm, dsfont, xspace, cancel, mathtools). The delivered numbering is the unit's own. Assets: no retained span contains a figure environment, a graphics-inclusion command, a PDF-page inclusion, a table or a TikZ picture; no image file is copied into this package, the package declares no asset dependency and no image file is redistributed with it. Licence: version arXiv:2608.12087v1 of this article is distributed under the Creative Commons Attribution 4.0 International licence, as recorded in the arXiv licence metadata for this exact version and shown on the abstract page https://arxiv.org/abs/2608.12087v1. A full rights notice is delivered at licenses/arxiv-2608.12087-CC-BY-4.0.txt. Characters: every retained excerpt is pure ASCII, so the delivered engine, XeLaTeX with its default Latin Modern text font, reports no missing character.
\begin{theorem}
\label{thm:variable-m-strongly-np-complete}
When \(m\) is part of the input, the decision problem
\[
        P,S1\mid s_j=t_j=1\mid C_{\max}
\]
is strongly NP-complete.
\end{theorem}

\begin{proof}
Membership in NP is immediate: a certificate specifies a machine assignment and
integer loading time for each job, and feasibility together with the bound
\(C_{\max}\le T\) can be verified in polynomial time.


The reduction is from \textsc{3-Partition}. An instance of \textsc{3-Partition} is a
set of \(3m\) positive integers
\[
        a_1,\ldots,a_{3m}
\]
such that
\[
        \sum_{i=1}^{3m}a_i=mB,
        \qquad
        \frac{B}{4}<a_i<\frac{B}{2}
        \quad (i=1,\ldots,3m).
\]
The question is whether the items can be partitioned into \(m\) disjoint triples
whose sums are all equal to \(B\). We may assume \(m\ge2\).

From this instance, we construct an instance of the loading-unloading scheduling
problem with \(m\) machines and \(3m\) jobs. For every item \(a_i\), create one
job \(J_i\) with execution length
\[
        e_i=2m a_i .
\]
Since \(e_i=p_i+2\), set
\[
        p_i=2m a_i-2 .
\]
As \(m\ge2\) and \(a_i\ge1\), these processing times are positive integers.
Finally, set the target makespan to
\[
        T=2mB+2m-2 .
\]

We show that the \textsc{3-Partition} instance is feasible if and only if the
constructed scheduling instance admits a schedule with makespan at most \(T\).

First, suppose that the items can be partitioned into triples
\[
        A_1,\ldots,A_m,
        \qquad
        \sum_{a_i\in A_j}a_i=B
        \quad (j=1,\ldots,m).
\]
For each machine \(j\), schedule the three jobs corresponding to \(A_j\)
consecutively, starting at time
\[
        L_j=2(j-1).
\]
If the three jobs on machine \(j\) are denoted by
\(J_{j,1},J_{j,2},J_{j,3}\), set
\[
        S_{j,1}=L_j,
        \qquad
        S_{j,h+1}=S_{j,h}+e_{j,h},
        \quad h=1,2.
\]
Thus, the jobs on each machine are processed without overlap. The total execution
length assigned to machine \(j\) is
\[
        \sum_{a_i\in A_j}2m a_i=2mB.
\]
Hence machine \(j\) completes at
\[
        L_j+2mB=2(j-1)+2mB.
\]
The largest completion time occurs on machine \(m\), and it is
\[
        2(m-1)+2mB=2mB+2m-2=T.
\]

It remains only to verify the server constraint. Every execution length is a
multiple of \(2m\). Hence all loading times on machine \(j\) are congruent to
\[
        2(j-1) \pmod{2m}.
\]
The unloading time of a job loaded at time \(S\) is
\[
        U=S+e_i-1.
\]
Therefore, all unloading times on machine \(j\) are congruent to
\[
        2(j-1)-1 \pmod{2m}.
\]
Thus, machine \(j\) uses one residue class modulo \(2m\) for its loadings and one
different residue class for its unloadings. Over all \(j=1,\ldots,m\), the
loading residues are
\[
        0,2,4,\ldots,2m-2,
\]
and the unloading residues are
\[
        2m-1,1,3,\ldots,2m-3.
\]
Together these are exactly the \(2m\) residue classes modulo \(2m\). Consequently
no two server operations coincide. This excludes loading-loading,
unloading-unloading, and loading-unloading collisions. The constructed schedule
is therefore feasible and has the makespan \(T\).

Conversely, suppose that the constructed scheduling instance has a feasible
schedule with
\[
        C_{\max}\le T.
\]
We prove that the original \textsc{3-Partition} instance has a feasible
partition.

No machine can process four jobs. Indeed, for any four jobs corresponding to
items \(a_{i_1},a_{i_2},a_{i_3},a_{i_4}\), we have
\[
        a_{i_1}+a_{i_2}+a_{i_3}+a_{i_4}>B.
\]
Since the item sizes are integral, this implies
\[
        a_{i_1}+a_{i_2}+a_{i_3}+a_{i_4}\ge B+1.
\]
The total execution length of the corresponding four jobs is therefore at least
\[
        2m(B+1)=2mB+2m.
\]
However, 
\[
        2mB+2m>T=2mB+2m-2.
\]
Thus, four jobs cannot be assigned to the same machine in any schedule with a 
makespan at most \(T\).

There are \(3m\) jobs and \(m\) machines. Since every machine processes at most
three jobs, every machine must process exactly three jobs.

Next, consider the total execution length assigned to any one machine. Each job
has an execution length divisible by \(2m\), and hence each machine load is a
multiple of \(2m\). No machine can have a load greater than \(2mB\). If a machine
load were larger than \(2mB\), then, being a multiple of \(2m\), it would be at
least
\[
        2m(B+1)=2mB+2m>T,
\]
which is impossible.

The total execution length of all jobs is
\[
        \sum_{i=1}^{3m} e_i
        =
        \sum_{i=1}^{3m}2m a_i
        =
        2m\cdot mB .
\]
Since there are \(m\) machines and each machine has a load at most \(2mB\), every
machine must have a load exactly equal to \(2mB\).

Therefore, each machine processes exactly three jobs, and the three jobs assigned
to any machine have total execution length \(2mB\). If these jobs correspond to
items \(a_i,a_j,a_k\), then
\[
        2m(a_i+a_j+a_k)=2mB,
\]
and hence
\[
        a_i+a_j+a_k=B.
\]
Therefore, the triples induced by the \(m\) machine assignments form a feasible
3-partition.

We have shown that the constructed scheduling instance has a schedule with a 
makespan at most \(T\) if and only if the original \textsc{3-Partition} instance
is feasible.

The transformation is polynomial, and all numerical values created by the
reduction are polynomially bounded in the numerical data of the
\textsc{3-Partition} instance. Since \textsc{3-Partition} is strongly
NP-complete, the reduction proves strong NP-hardness. Together with membership
in NP, the decision problem is strongly NP-complete.
\end{proof}

\end{document}
